\documentclass{article}

\usepackage{PRIMEarxiv}
\usepackage[utf8]{inputenc}
\usepackage[T1]{fontenc}
\usepackage{hyperref}
\usepackage{url}
\usepackage{booktabs}
\usepackage{amsfonts}
\usepackage{amsmath}
\usepackage{amssymb}
\usepackage{nicefrac}
\usepackage{microtype}
\usepackage{fancyhdr}
\usepackage{graphicx}

\graphicspath{{media/}}
\hypersetup{hidelinks}
\urlstyle{same}

\pagestyle{fancy}
\thispagestyle{empty}
\rhead{\textit{}}
\fancyhead[LO]{Fairness, efficiency, and strategic behaviour in shared GPU allocation}

\title{Fairness, efficiency, and strategic behaviour\\in shared GPU allocation}

\author{
  Aayush Kuloor \quad Harsh Dhru \quad Raj Modi \quad Kalp Shah \\
  Indian Institute of Technology Gandhinagar \\
  \texttt{\{aayush.kuloor, harsh.dhru, raj.modi, kalp.shah\}@iitgn.ac.in}
}

\begin{document}
\maketitle

\begin{abstract}
Shared GPU clusters have to divide scarce compute among users whose job values
are private and whose reports may be strategic. We study this problem in a
repeated, unit-demand allocation model. The evaluation measures allocative
efficiency, long-run fairness, and resistance to manipulation. A common
simulator compares five mechanisms: random allocation, round-robin, greedy
allocation by reported value, a history-penalised scoring rule, and a $k$-unit
Vickrey auction. Reports are bounded. Payment and non-payment mechanisms use
separate models, and all comparisons use paired valuation traces. The
experiments vary scarcity, the fairness penalty, the fraction and strategy of
misreporting users, valuation heterogeneity, temporal persistence, and
population size. We will estimate the resulting frontier over welfare,
fairness, and manipulation gain. The output will be a reproducible stress test
of simple allocation rules, with evidence about when a modest fairness
correction helps, fails, or introduces new strategic incentives.
\end{abstract}

\keywords{GPU allocation \and fair resource allocation \and mechanism design
  \and strategic behaviour \and multi-agent systems}

% ============================================================
\section{Introduction}
\label{sec:intro}
% ============================================================

Modern machine-learning workloads compete for expensive shared accelerators.
The scheduling rule affects cluster efficiency and the time each user waits.
Existing systems handle priorities, heterogeneous accelerators, placement, and
long-running jobs \cite{gu2019tiresias,narayanan2020heterogeneity}. Fair access
still matters. Repeatedly serving the jobs with the largest apparent benefit
can starve other users. A rigid equal-time rule can also waste compute when job
values differ substantially.

Private information creates another difficulty. If allocation depends on a
self-reported priority, a user may inflate the report to receive service. The
scheduler sees the claim, not the user's true value. A rule that is efficient
with truthful reports can therefore perform poorly when users act
strategically. The analysis draws on online resource allocation, fair division,
and mechanism design
\cite{nisan2007algorithmic,ghodsi2011dominant,bertsimas2011price}.

Several GPU schedulers provide useful points of comparison. Tiresias uses
attained service to reduce job completion time. Gavel extends scheduling
policies to heterogeneous accelerators. Themis combines finish-time fairness
with a round-by-round partial-allocation auction designed to encourage truthful
bids. Shockwave uses market-based planning for jobs that adapt over time
\cite{gu2019tiresias,narayanan2020heterogeneity,mahajan2020themis,
zheng2023shockwave}. This project studies a narrower question. It uses a
stylised unit-demand model to separate the effects of scarcity, allocation
history, and bounded misreporting while holding information and random inputs
constant. It does not claim a new production scheduler or a general
impossibility theorem.

\paragraph{Research question.}
\begin{quote}
\emph{How do simple repeated GPU-allocation mechanisms trade off true social
welfare, long-run fairness, and empirical resistance to bounded strategic
reporting as scarcity, user heterogeneity, and temporal demand patterns vary?}
\end{quote}

The experiments address three subsidiary questions:
\begin{enumerate}
  \item Does a moderate history penalty reduce starvation without a large loss
        of true welfare?
  \item Does increasing the fairness penalty change the benefit available from
        strategically timing or inflating reports?
  \item Do conclusions based on equal allocation counts survive when users
        have unequal value distributions or persistent demand?
\end{enumerate}

\paragraph{Research contribution.}
The project places allocation and payments in one formal model. It measures
manipulation gain with paired counterfactual simulations and evaluates welfare,
fairness, and manipulation together. Validation tests, confidence intervals,
saved seeds, and common workloads make the experiments reproducible. The study
is an analysis of mechanisms under controlled workloads, not a systems
implementation.

% ============================================================
\section{Preliminaries}
\label{sec:prelim}
% ============================================================

\paragraph{Agents, resources, and time.}
Let $[n]=\{1,\ldots,n\}$ denote the unit-demand users. In each of $T$
discrete rounds, they compete for $k<n$ identical GPUs. A user can receive at
most one GPU in a round. Since all values are non-negative, the base model
allocates all $k$ GPUs. A seeded random generator breaks ties uniformly.

\paragraph{Valuations and reports.}
User $i$ draws a private value $v_{i,t}\in[0,v_{\max}]$ from distribution
$\mathcal D_i$ for receiving one GPU in round $t$. The user's report is
$\hat v_{i,t}\in[0,v_{\max}]$. The common bound on values and reports rules
out the degenerate case in which an arbitrarily large bid controls the
allocation. Unless stated otherwise, $v_{\max}=1$.

\paragraph{Allocations, payments, and histories.}
An allocation $\mathbf x_t\in\{0,1\}^n$ satisfies
$\sum_i x_{i,t}=k$. Its payment vector
$\mathbf p_t\in\mathbb R_{\ge 0}^n$ has $p_{i,t}=0$ whenever
$x_{i,t}=0$. For non-monetary mechanisms, $\mathbf p_t=\mathbf0$. The public
history before round $t$ is
$h_t=(\mathbf x_1,\ldots,\mathbf x_{t-1})$, and
\[
  a_i(t)=\sum_{\tau=1}^{t-1}x_{i,\tau}
\]
is user $i$'s cumulative allocation. Let $q_i(t)$ be the number of
consecutive rounds since user $i$ last received a GPU.

\paragraph{Utility and welfare.}
Under quasi-linear utility, user $i$ receives
\[
  u_{i,t}=v_{i,t}x_{i,t}-p_{i,t},
  \qquad U_i=\sum_{t=1}^T u_{i,t}.
\]
The gross allocation benefit is $G_i=\sum_t v_{i,t}x_{i,t}$. Payments count as
transfers when calculating social welfare:
\[
  W=\sum_{i=1}^nG_i
   =\sum_{t=1}^T\sum_{i=1}^n v_{i,t}x_{i,t}.
\]
In round $t$, let $S_t^*$ contain the $k$ users with the highest true values.
The truthful allocation oracle obtains
\[
  W^*=\sum_{t=1}^T\sum_{i\in S_t^*}v_{i,t},
  \qquad \mathrm{WR}=W/W^*.
\]

\paragraph{Allocation and benefit fairness.}
Let $A_i=a_i(T+1)$. Jain's allocation-fairness index is
\cite{jain1984quantitative}
\[
  J_A=\frac{(\sum_i A_i)^2}{n\sum_i A_i^2}.
\]
When the distributions $\mathcal D_i$ differ, equal access need not yield
equal benefit. We also report Jain's index $J_B$ on normalised gross benefits:
\[
  B_i=\frac{G_i}{T\mu_i},\qquad
  \mu_i=\mathbb E_{v\sim\mathcal D_i}[v],\qquad
  J_B=\frac{(\sum_iB_i)^2}{n\sum_iB_i^2}.
\]
Nash gross welfare is
$\mathrm{NSW}=\sum_i\log(G_i+\varepsilon)$ with
$\varepsilon=10^{-8}$.

\paragraph{Waiting and starvation.}
The maximum wait is $Q_{\max}=\max_{i,t}q_i(t)$. For
$\Delta=2\lceil n/k\rceil$, twice the service cycle of round-robin, the
starvation rate is
\[
  \mathrm{SR}_{\Delta}
  =\frac{1}{nT}\sum_{i=1}^n\sum_{t=1}^T
    \mathbb I\{q_i(t)>\Delta\}.
\]
We report the $95$th percentile of all waiting times with this threshold-based
measure.

\paragraph{Paired strategic effects.}
Let $\omega$ contain a fixed valuation trace, tie-breaking seed, and set of
strategic users. To evaluate manipulation gain, we change only the focal user's
strategy and hold $\omega$ and every other policy fixed:
\[
  M_i(\sigma_i;\omega)
  =U_i(\sigma_i,\sigma_{-i};\omega)
   -U_i(\text{truthful},\sigma_{-i};\omega).
\]
The reported statistics are the mean $\overline M$, the maximum $M_{\max}$,
and the fraction of strategic users with $M_i>0$. For a strategic profile, the
welfare loss is
\[
  \mathrm{PoS}
  =\frac{W_{\mathrm{truthful}}-W_{\mathrm{strategic}}}
         {W_{\mathrm{truthful}}},
\]
Both simulations use the same $\omega$. For a fair mechanism $F$, the
normalised welfare loss with truthful reports is
\[
  \mathrm{PoF}(F)=\frac{W^*-W_F}{W^*}.
\]

\paragraph{Incentive terminology.}
A one-round direct mechanism is dominant-strategy incentive compatible (DSIC)
when truthful reporting weakly maximises each user's utility for every profile
of other reports \cite{nisan2007algorithmic}. Reports cannot affect a
report-invariant rule, which makes such a rule weakly strategyproof. For
history-dependent non-payment rules, the experiments measure vulnerability to
specified attacks. Simulation cannot establish dynamic strategyproofness.

% ============================================================
\section{Problem formulation}
\label{sec:formulation}
% ============================================================

Let $\rho\in[0,1]$, and let
$\mathcal S\subseteq[n]$, $|\mathcal S|=\lfloor\rho n\rfloor$, be the
strategic set. A reporting policy may condition on the current value and
public history:
\[
  \sigma_i:[0,v_{\max}]\times\mathcal H_t
  \longrightarrow[0,v_{\max}],
  \qquad
  \hat v_{i,t}=\sigma_i(v_{i,t},h_t).
\]
A repeated mechanism is a sequence of mappings
\[
  \mathcal M_t:
  [0,v_{\max}]^n\times\mathcal H_t
  \longrightarrow
  \mathcal X_k\times\mathbb R_{\ge0}^n,
  \qquad
  \mathcal X_k=
  \left\{\mathbf x\in\{0,1\}^n:\sum_i x_i=k\right\},
\]
which returns an allocation $\mathbf x_t$ and a payment vector $\mathbf p_t$.
The same model therefore covers the non-payment rules and the auction.

For each workload condition, we estimate the achievable set
\[
  \mathcal F=
  \left\{
  (\mathrm{WR},J_A,J_B,\mathrm{SR}_{\Delta},\overline M):
  \mathcal M\in\mathfrak M
  \right\},
\]
where $\mathfrak M$ contains the five candidate mechanisms and the parameter
settings described below. One setting Pareto dominates another when it has no
lower welfare or fairness, no higher starvation or manipulation gain, and a
strict improvement on at least one measure. Keeping the measures separate
avoids an arbitrary scalar weighting of these normative objectives.

The base environment is a repeated allocation model, not a production
scheduler. It assumes identical GPUs, one-round service, at most one GPU per
user, and additive values. It omits job arrivals, gang scheduling, placement,
preemption, and multi-GPU speedup. These restrictions allow controlled causal
comparisons between mechanisms. A separate experiment introduces temporal
persistence to test whether the conclusions depend on independent values.

% ============================================================
\section{Proposed solutions}
\label{sec:solutions}
% ============================================================

\subsection{Allocation mechanisms}

\begin{table}[ht]
\caption{Properties evaluated or guaranteed by construction.}
\label{tab:mechanisms}
\centering
\small
\begin{tabular}{lcccc}
\toprule
Mechanism & Reports & History & Payments & Formal property \\
\midrule
M1 Random      & No  & No  & No  & Report-invariant \\
M2 Round-robin & No  & Yes & No  & Report-invariant \\
M3 Greedy      & Yes & No  & No  & Truthful welfare oracle only \\
M4 Score       & Yes & Yes & No  & Empirically stress-tested \\
M5 Vickrey     & Yes & No  & Yes & Per-round DSIC \\
\bottomrule
\end{tabular}
\end{table}

\paragraph{M1 Random.}
M1 chooses $k$ users uniformly without replacement. It ignores reports and
history. It is report-invariant and equalises expected allocation counts.

\paragraph{M2 Round-robin.}
M2 starts from a seeded queue and serves the next $k$ users in each round. It
is report-invariant and bounds waiting time, but ignores current job value.

\paragraph{M3 Greedy reported value.}
Allocate to the $k$ largest reports:
\[
  \mathbf x_t\in
  \operatorname*{arg\,max}_{\mathbf x\in\mathcal X_k}
  \sum_{i=1}^n\hat v_{i,t}x_i.
\]
When reports equal values, M3 attains the truthful welfare oracle. Increasing a
report can only improve the user's one-round allocation probability.

\paragraph{M4 History-penalised score.}
For $\lambda\ge0$, define
\[
  s_{i,t}=\frac{\hat v_{i,t}}{(1+a_i(t))^{\lambda}}
\]
and allocate to the $k$ largest scores. M4 equals M3 at $\lambda=0$. As
$\lambda$ grows, cumulative service has more influence relative to reported
value. The rule approaches least-served-first, with reports used to break ties,
but never becomes literal round-robin. Sweeping $\lambda$ measures the fairness
benefit and strategic cost of the correction.

\paragraph{M5 $k$-unit Vickrey auction.}
Allocate to the $k$ highest bids and charge each winner the
$(k+1)$st-highest bid:
\[
  p_{i,t}=x_{i,t}\hat v_{(k+1),t}.
\]
With unit demand, quasi-linear utility, and no intertemporal budget, this is the
VCG outcome for $k$ identical units and is DSIC in each round
\cite{vickrey1961counterspeculation,nisan2007algorithmic}. The model treats
payments as utility-equivalent transfers. They are not reusable virtual
credits. Real transfers are required for this truthfulness benchmark, which we
use only as a mechanism-design baseline and do not propose as a deployment
policy.

\subsection{Strategic policies}

Each strategic user follows one of four bounded policies:
\begin{enumerate}
  \item \textbf{Truthful:}
        $\hat v_{i,t}=v_{i,t}$.
  \item \textbf{Capped exaggeration:}
        $\hat v_{i,t}=\min\{c v_{i,t},v_{\max}\}$ for
        $c\in\{1.25,1.5,2\}$.
  \item \textbf{Maximum claim:}
        $\hat v_{i,t}=v_{\max}$, a worst-case priority-inflation baseline.
  \item \textbf{Finite-horizon rollout attack:}
        for a diagnostic setting with $n=10$, evaluate reports on
        $\mathcal G=\{0,0.1,\ldots,1\}$ using $100$ Monte Carlo rollouts over
        the next $H=5$ rounds, and select the report with highest estimated
        cumulative utility under fixed opponent policies.
\end{enumerate}
The rollout policy approximates a unilateral attack. It is neither an
equilibrium computation nor a proof of manipulability. We report it separately
because a finite search heuristic is not a true best response.

\subsection{Experimental protocol}

Unless a factor is being varied, the base configuration is
\[
  n=50,\quad k=10,\quad T=1000,\quad
  \mathcal D_i=\mathrm{Uniform}(0,1),\quad \rho=0.25.
\]
Each condition uses $30$ independent master seeds. For a given seed, every
mechanism receives the same valuation tensor, strategic-user set, and
pre-generated tie-breaking randomness. Reported results will include seed
means and paired $95\%$ bootstrap confidence intervals based on $10{,}000$
resamples. We will save raw per-seed results, configuration files, and the
master seed list. Experiment E3 uses a factorial design; all other experiments
change one factor at a time.

\paragraph{E0 validation and sanity checks.}
Small discrete instances allow exhaustive checks of capacity and payment
invariants, report-invariance of M1--M2, equivalence of M3 and M4 at
$\lambda=0$, the round-robin wait bound, welfare-oracle calculations, and
one-round truthfulness of M5 for every report in a finite grid. These checks
identify implementation errors before the main experiments run.

\paragraph{E1 resource scarcity.}
Vary $k/n\in\{0.1,0.2,0.4,0.6,0.8\}$ and compare WR, $J_A$,
$Q_{\max}$, the $95$th-percentile wait, and $\mathrm{SR}_{\Delta}$ under
truthful reports.

\paragraph{E2 fairness and strategic-vulnerability frontier.}
Sweep
$\lambda\in\{0,0.05,0.1,0.25,0.5,1,2,5\}$. Plot WR against $J_A$ and
colour each point by $\overline M$ under capped exaggeration. Report all
non-dominated settings. No single metric will determine the selected value of
$\lambda$.

\paragraph{E3 strategic population and attack type.}
Use a factorial design over
$\rho\in\{0,0.1,0.25,0.5,1\}$ and the three scalable policies (truthful,
capped exaggeration with $c=2$, and maximum claim). Measure PoS,
$\overline M$, $M_{\max}$, the fraction benefiting from manipulation, WR,
  and fairness. Apply the rollout attack to selected settings for M3--M5 where
  it is most likely to separate their behaviour.

\paragraph{E4 heterogeneous users.}
Compare the homogeneous base case with an equal-sized two-group population:
$\mathcal D_L=\mathrm{Beta}(2,5)$ and
  $\mathcal D_H=\mathrm{Beta}(5,2)$. Both remain in $[0,1]$ but have different
  means. Report $J_A$ and $J_B$ together to identify cases in which equal
  service and equal normalised benefit diverge.

\paragraph{E5 temporally persistent demand.}
Replace i.i.d. values with
\[
  v_{i,t}=\alpha v_{i,t-1}+(1-\alpha)\epsilon_{i,t},
  \qquad \epsilon_{i,t}\sim\mathrm{Uniform}(0,1),
\]
for $\alpha\in\{0,0.5,0.9\}$. The experiment checks whether the conclusions
hold when high-value demand persists across rounds, a feature absent from the
i.i.d. model.

\paragraph{E6 scalability.}
Vary $n\in\{10,25,50,100,250,500\}$ with $k/n=0.2$ and $T=1000$.
Report wall-clock time per simulated round and peak memory. This experiment
excludes the rollout attack, which solves a separate and more expensive
diagnostic problem.

\subsection{Testable hypotheses}

\begin{enumerate}
  \item Moderate $\lambda$ will reduce tail waiting and starvation before it
        causes a comparable loss in WR.
  \item Increasing $\lambda$ will improve allocation equality but may increase
        the value of strategically timing reports because current service
        changes future scores.
  \item Round-robin will score well on $J_A$ but worse on WR and, under
        heterogeneous values, need not dominate on $J_B$.
  \item M5 will have negligible positive manipulation gain in the one-round
        validation and repeated independent setting, whereas M3 and M4 will
        remain vulnerable to bounded over-reporting.
\end{enumerate}
The project can succeed even if the evidence rejects these hypotheses.

\subsection{Anticipated limitations}

The model omits job duration, multi-GPU jobs, placement, and preemption. Its
findings apply to repeated unit-demand allocation, not to a production cluster.
Synthetic values and utility-equivalent payments cannot establish external
validity. Jain's index captures one notion of equality, so the analysis also
reports waiting and normalised-benefit metrics. Fixed attack policies do not
characterise a multi-agent equilibrium, and the rollout attack is a
finite-horizon unilateral search. The Vickrey result also depends on real
quasi-linear transfers. Replacing payments with reusable credits would require
a separate dynamic budget model and would remove the stated DSIC guarantee.

\subsection{Success criteria and reproducibility}

The project is \emph{complete} when M1--M5 pass E0, E1--E4 run for every
prespecified seed, paired intervals and raw results are produced, and one
command can regenerate every main figure. E5--E6 and the rollout attack are
stretch deliverables. The evidence must answer the research questions
reliably, even when it contradicts the hypotheses.

% ============================================================
\section{Contribution}
\label{sec:contributions}
% ============================================================

All members will review the formal model, experimental design, results, and
final report. Each member has a primary area of responsibility.

\paragraph{Aayush Kuloor: simulator and reproducibility.}
Aayush will implement the valuation processes, histories, paired random
streams, configuration loading, result storage, and end-to-end experiment
runner. He will maintain the master seeds and reproducibility instructions.

\paragraph{Harsh Dhru: mechanisms and formal validation.}
Harsh will implement M1--M5, the allocation and payment interfaces, seeded
tie-breaking, and the exhaustive E0 tests for mechanism invariants and the
Vickrey benchmark.

\paragraph{Raj Modi: strategic policies and metrics.}
Raj will implement the bounded attacks and rollout diagnostic. He will also
implement WR, $J_A$, $J_B$, waiting, starvation, PoF, PoS, and paired
manipulation metrics, then validate them on hand-computed examples.

\paragraph{Kalp Shah: experiments and statistical analysis.}
Kalp will implement the parameter sweeps, paired bootstrap intervals,
Pareto-dominance analysis, figures, and result tables. He will document
failures, sensitivity checks, and positive findings.

Each owner will provide unit tests and a short methods description for their
component. Another member will review that component's code.

\bibliographystyle{unsrt}
\bibliography{references}

\end{document}
